\begin{helper}
Let \(g:\mathbb R^3\to\mathbb R\) be continuous, written
\(g(x,y,z)\), and suppose \(g(x,y,z)\ge 0\) on the simplex
\(0\le x\le y\le z\le 1\).  Then the iterated nonnegative integral of
\(\operatorname{ofReal}(g)\) over this simplex is bounded by the
\(\operatorname{ofReal}\) of the corresponding iterated real interval
integral:
\[
\int_{\mathbb R}^{+}\int_{\mathbb R}^{+}\int_{\mathbb R}^{+}
\mathbf 1_{\{0\le x\le y\le z\le 1\}}\operatorname{ofReal}(g(x,y,z))
\,dz\,dy\,dx
\le
\operatorname{ofReal}\left(\int_0^1\int_0^z\int_0^y
g(x,y,z)\,dx\,dy\,dz\right).
\]
\end{helper}
\begin{proof}
Put
\[
S=\{(x,y,z)\in\mathbb R^3:0\le x\le y\le z\le 1\}
\]
and
\[
H(x,y,z)=\mathbf 1_S(x,y,z)\operatorname{ofReal}(g(x,y,z)).
\]
The set \(S\) is Borel measurable: it is the intersection of the closed
half-spaces cut out by the continuous functions \(x\), \(y-x\), \(z-y\), and
\(1-z\).  Since \(g\) is continuous, the function
\((x,y,z)\mapsto \operatorname{ofReal}(g(x,y,z))\) is Borel measurable as an
\([0,\infty]\)-valued function, and therefore the indicator product \(H\) is
measurable and nonnegative.  The Tonelli/order-swap and restriction lemma
\(\noderef{SamplingSimplexIndicatorLIntegralRestriction}\), applied to this
continuous \(g\), gives the exact equality
\[
\int_{\mathbb R}^{+}\int_{\mathbb R}^{+}\int_{\mathbb R}^{+}H(x,y,z)
\,dz\,dy\,dx
=
\int_{(0,1]}^{+}\int_{(0,z]}^{+}\int_{(0,y]}^{+}
\operatorname{ofReal}(g(x,y,z))\,dx\,dy\,dz .
\]

It remains to convert the nested nonnegative interval integral into
\(\operatorname{ofReal}\) of the nested real interval integral.  The
one-dimensional conversion fact is
\(\noderef{SamplingIntervalLIntegralOfRealConversion}\): on an oriented
interval \([a,b]\) with \(a\le b\), if a real function is interval-integrable
and nonnegative on \((a,b]\), then the lintegral of \(\operatorname{ofReal}\)
over \((a,b]\) equals \(\operatorname{ofReal}\) of the real interval integral.
The regularity and nonnegativity hypotheses needed for the three applications
are supplied by
\(\noderef{SamplingSimplexPartialIntegralRegularity}\), applied to the present
continuous \(g\) and its assumed simplex nonnegativity.

First fix \(y,z\) with \(0\le y\le z\le 1\).  The function
\(x\mapsto g(x,y,z)\) is interval-integrable on \([0,y]\) and nonnegative on
\((0,y]\) by the first clause of
\(\noderef{SamplingSimplexPartialIntegralRegularity}\).  Since \(0\le y\),
\(\noderef{SamplingIntervalLIntegralOfRealConversion}\) gives
\[
\int_{0}^{y,+}\operatorname{ofReal}(g(x,y,z))\,dx
=
\operatorname{ofReal}\left(\int_0^y g(x,y,z)\,dx\right).
\]

Second fix \(z\in[0,1]\) and define
\[
F_z(y)=\int_0^y g(x,y,z)\,dx .
\]
The second clause of
\(\noderef{SamplingSimplexPartialIntegralRegularity}\) says that \(F_z\) is
continuous on \([0,z]\), interval-integrable on \([0,z]\), and nonnegative on
\((0,z]\).  Since \(0\le z\),
\(\noderef{SamplingIntervalLIntegralOfRealConversion}\) gives
\[
\int_0^{z,+}\operatorname{ofReal}\left(\int_0^y g(x,y,z)\,dx\right)\,dy
=
\operatorname{ofReal}\left(\int_0^z\int_0^y g(x,y,z)\,dx\,dy\right).
\]

Third define
\[
G(z)=\int_0^z\int_0^y g(x,y,z)\,dx\,dy .
\]
The third clause of \(\noderef{SamplingSimplexPartialIntegralRegularity}\)
says that \(G\) is continuous on \([0,1]\), interval-integrable on \([0,1]\),
and nonnegative on \((0,1]\).  Since \(0\le1\), a third application of
\(\noderef{SamplingIntervalLIntegralOfRealConversion}\) yields
\[
\int_0^{1,+}\operatorname{ofReal}\left(\int_0^z\int_0^y
g(x,y,z)\,dx\,dy\right)\,dz
=
\operatorname{ofReal}\left(\int_0^1\int_0^z\int_0^y
g(x,y,z)\,dx\,dy\,dz\right).
\]
Combining the Tonelli order swap, the restriction of the indicator integral to
the simplex bounds, and the three \(\operatorname{ofReal}\)-conversion steps
shows that the left side of the statement is equal to the right side.  The
claimed inequality follows immediately.
\end{proof}
